% Generated by Listas/make_statements.py from pratica-lecture-04.md, item 1. Do not edit.
\textbf{Source.} Tadelis (2013), \emph{Game Theory: An Introduction}, Exercise 8.8 ``Entry Deterrence 1'', pp.~171--172 {[}PDF 190--191{]}. Reproduced verbatim. Answer in the source: none in the book (the preface, p.~xiii, says about half are solved on Princeton UP's website; not checked).



\begin{quote}

\textbf{Entry Deterrence 1:} NSG Corporation is considering entry into the local phone market in the San Francisco Bay Area. The incumbent, S\&P Corporation, predicts that a price war will result if NSG enters. If NSG stays out, S\&P earns monopoly profits valued at \$10 million (net present value of profits) while NSG earns zero. If NSG enters, it must incur irreversible entry costs of \$2 million. If there is a price war, each firm earns \$1 million. S\&P always has the option of accommodating entry (i.e., not starting a price war). In such a case both firms earn \$4 million (net present value). Suppose that the timing is such that NSG first has to choose whether or not to enter the market. Then S\&P decides whether to ``accommodate entry'' or ``engage in a price war.'' What is the subgame-perfect equilibrium outcome to this sequential game? Set up a game tree.

\end{quote}

